\documentclass{standalone}
\usepackage{tikzplot}

\begin{document}

% begin_snippet
\begin{tikzpicture}
  \tikzmath{
    \A{1,1} = 0; \A{2,1} = 2; \A{3,1} = 1;
    \B{1,1} = 3; \B{2,1} = 0; \B{3,1} = 1;
    \C{1,1} = -4; \C{2,1} = -1; \C{3,1} = 1;
  }
  \tikzset{
    % 过 A,B,C 三点的圆
    conics/define/circle={\A,\B,\C,\Ca},
    % 以 A 为圆心 AB 为半径的圆
    conics/define/circle={\B,\A,\Cb},
    % 以 A 为圆心 AC 为半径的圆
    conics/distance={\A,\B,\d},
    conics/define/circle/center-radius={\A,\d,\Cc},
  }
  \coordinate (A) at (\A{1,1},\A{2,1});
  \coordinate (B) at (\B{1,1},\B{2,1});
  \coordinate (C) at (\C{1,1},\C{2,1});
  \axes[grid] (-5:5,-5:5);
  \conic[thick,draw=teal] (\Ca);
  \conic[thick,draw=violet] (\Cb);
  \conic[thick,draw=magenta] (\Cc);

  \foreach \p/\placement in {A/above,B/above,C/below right}{
    \fill[red] (\p) circle (1.5pt);
    \draw (\p) node[\placement] {$\p$};
  }
\end{tikzpicture}
% end_snippet

\end{document}
